\documentclass{article}
\usepackage{amsmath, amssymb, bm, array, booktabs, geometry}
\geometry{margin=1in}

\title{Mathematical Derivations of Fast Topological Data Analysis\\
	(Ellipse Fitting via Takens Embedding)}
\author{}
\date{}

\begin{document}
	\maketitle
	\tableofcontents
	\newpage
	

	\section{Spectral Analysis and TDA Parameter Selection}

	
	\textbf{Goal:} Extract two key frequencies from the signal ---
	$F_{\max}$ (dominant oscillation) and $F_{\min}$ (lowest
	significant frequency) --- which control the time delay
	and window size used in all subsequent steps.
	Transforming the signal from the time domain into
	the frequency domain using the DFT, then locate peaks in the
	resulting amplitude spectrum.
	
	\subsection{Discrete Fourier Transform}
	
	The DFT decomposes the time-domain signal $x(n)$ into its
	sinusoidal frequency components.
	Each coefficient $Y(k)$ tells us how much energy the signal
	has at frequency $f(k) = k \cdot F_s/N$ Hz.
	Given a discrete time series $x(n)$ of length $N$:
	
	\begin{equation}
		Y(k) = \sum_{n=0}^{N-1} x(n)\, e^{-j2\pi kn/N},
		\qquad k = 0, 1, \ldots, N-1
	\end{equation}
	
	Because the input is real-valued, the DFT is symmetric and
	only the first half carries unique information.
	Dividing by $N$ normalises the amplitude so it is independent
	of signal length, and doubling the interior bins preserves
	total energy in the single-sided representation:
	
	\begin{equation}
		|Y_{\text{single}}(k)| =
		\begin{cases}
			\dfrac{|Y(k)|}{N}    & k = 0 \text{ (DC)} \\[6pt]
			\dfrac{2\,|Y(k)|}{N} & k = 1, \ldots, \lfloor N/2 \rfloor - 1 \\[6pt]
			\dfrac{|Y(k)|}{N}    & k = N/2 \text{ (Nyquist, even } N \text{ only)}
		\end{cases}
	\end{equation}
	
	The corresponding frequency axis (Hz) is:
	
	\begin{equation}
		f(k) = k \cdot \frac{F_s}{N},
		\qquad k = 0, 1, \ldots, \lfloor N/2 \rfloor
	\end{equation}
	
	where $F_s = 1/\Delta t$ is the sampling frequency.
	
	\subsection{Peak Frequency Detection}
	
	Two characteristic frequencies are extracted from the spectrum.
	$F_{\max}$ is the dominant oscillation frequency.
	$F_{\min}$ is the lowest frequency with a peak prominence above
	10\% of the maximum magnitude, ensuring noise bumps are ignored:
	
	\begin{align}
		F_{\max} &= \text{frequency at absolute spectral peak} \\[4pt]
		F_{\min} &= \text{lowest peak with prominence} > 10\%
		\text{ of } \max|Y_{\text{single}}|
	\end{align}
	
	\subsection{TDA Time Delay and Window Duration}
	
	\textbf{Goal:} Translate $F_{\max}$ and $F_{\min}$ into
	the two control parameters of the Takens embedding:
	the time delay $\tau$ and the window duration $T_w$.
	Applying the Razmarashooli method, which sets
	$\tau$ to one quarter of the dominant period and $T_w$
	to one full period of the lowest frequency.
	
	Following the Razmarashooli method, the time delay $\tau$ and
	window duration $T_w$ are defined as:
	
	\begin{align}
		\tau      &= \frac{0.25}{F_{\max}} \quad \text{(seconds)}
		\label{eq:tau} \\[4pt]
		\tau_s    &= \operatorname{round}\!\left(\frac{\tau}{\Delta t}\right)
		\quad \text{(samples)} \\[8pt]
		T_w       &= \frac{1}{F_{\min}} \quad \text{(seconds)}
		\label{eq:Tw} \\[4pt]
		N_w       &= \operatorname{round}\!\left(\frac{T_w}{\Delta t}\right)
		\quad \text{(samples per window)}
	\end{align}
	
	\textbf{Rationale for $\tau = T/4$:}
	For a sinusoid $x(t) = A\sin(2\pi F_{\max} t)$, a
	quarter-period delay produces a cosine:
	
	\begin{equation}
		x(t + \tau) = A\sin\!\left(2\pi F_{\max} t + \frac{\pi}{2}\right)
		= A\cos(2\pi F_{\max} t)
	\end{equation}
	
	Plotting $x(t)$ vs.\ $x(t+\tau)$ traces a perfect circle,
	maximising the 2-D spread of the point cloud and giving
	the best-conditioned ellipse fit.
	Any other delay angle produces an ellipse, but $\tau = T/4$
	is the unique delay that makes the two embedded dimensions
	statistically independent for a sinusoid.
	

	\section{Zero-Padding}
	
	
	\textbf{Goal:} Ensure the sliding window is valid from sample
	index $i=0$ without accessing samples that do not exist
	in the original signal.
	Prepending $N_{\text{pad}}$ zeros to the signal
	before the main loop begins.
	
	The earliest sample accessed by the window at position $i=0$
	is at conceptual index:
	
	\begin{equation}
		\ell = -(N_w - 1) - \tau_s < 0
	\end{equation}
	
	Since $\ell < 0$, samples before the signal start are required.
	The required number of prepended zeros is:
	
	\begin{equation}
		N_{\text{pad}} = |\ell| = N_w - 1 + \tau_s
	\end{equation}
	
	The zero-padded signal is defined as:
	
	\begin{equation}
		x_{\text{pad}}(n) =
		\begin{cases}
			0                  & n = 0, \ldots, N_{\text{pad}}-1
			\quad \text{(prepended zeros)} \\
			x(n - N_{\text{pad}}) & n = N_{\text{pad}}, \ldots,
			N_{\text{pad}}+N-1
			\quad \text{(original signal)}
		\end{cases}
	\end{equation}
	

	\section{Takens Delay Embedding}
	
	
	\textbf{Goal:} Convert the 1-D signal within each window
	into a 2-D point cloud that traces an ellipse in phase space.
    Pairing each sample with a delayed copy of itself,
	separated by $\tau_s$ samples.
	For a periodic signal this produces a closed elliptical loop;
	deviations from the loop reveal dynamic changes.
	
	For window starting at sample index $i$, the two-dimensional
	point cloud $\mathbf{P}$ is constructed as:
	
	\begin{equation}
		\mathbf{P} =
		\begin{bmatrix}
			x(i)                  & x(i + \tau_s) \\
			x(i+1)                & x(i+\tau_s+1) \\
			\vdots                & \vdots \\
			x(i+N_w-\tau_s-1)    & x(i+N_w-1)
		\end{bmatrix}
		\in \mathbb{R}^{(N_w - \tau_s) \times 2}
	\end{equation}
	
	Each row $k$ is a 2-D point in the reconstructed phase space:
	
	\begin{equation}
		\mathbf{p}_k = \bigl(x(i+k-1),\; x(i+k-1+\tau_s)\bigr),
		\qquad k = 1, \ldots, N_w - \tau_s
	\end{equation}
	
	Let $n = N_w - \tau_s$ denote the number of points per window.
	
	
	\section{General Conic Equation and Ellipse Constraint}
	
	
	\textbf{Goal:} Define the mathematical form of the curve we
	are fitting (an ellipse), and establish the algebraic
	constraint that guarantees the fitted curve is always
	an ellipse --- never a hyperbola or parabola.
    Useing the general degree-2 implicit equation and
	apply the Fitzgibbon constraint $4ac - b^2 = 1$.
	
	\subsection{General Conic}
	
	Any conic section in $\mathbb{R}^2$ satisfies the implicit equation:
	
	\begin{equation}
		\boxed{F(x,y) = ax^2 + bxy + cy^2 + dx + ey + f = 0}
		\label{eq:conic}
	\end{equation}
	
	with coefficient vector $\mathbf{a} = [a,\, b,\, c,\, d,\, e,\, f]^\top$.
	The coefficients $a$, $b$, $c$ control shape and tilt;
	$d$, $e$, $f$ control position and offset.
	
	\subsection{Ellipse Discriminant Condition}
	
	The conic type is determined by the discriminant $\Delta = b^2 - 4ac$:
	
	\begin{equation}
		\Delta =
		\begin{cases}
			< 0 & \Rightarrow \text{ellipse} \\
			= 0 & \Rightarrow \text{parabola} \\
			> 0 & \Rightarrow \text{hyperbola}
		\end{cases}
	\end{equation}
	
	Without a constraint, an unconstrained fit might return
	a hyperbola or parabola that also passes through the data.
	We must explicitly enforce $\Delta < 0$.
	
	\subsection{Fitzgibbon Algebraic Constraint}
	
	To guarantee an ellipse solution, Fitzgibbon et al.\ enforce:
	
	\begin{equation}
		\mathbf{a}^\top \mathbf{C}\, \mathbf{a} = 1,
		\qquad \text{where} \quad
		4ac - b^2 = 1
		\label{eq:constraint}
	\end{equation}
	
	This simultaneously fixes the scale ambiguity of $\mathbf{a}$
	(since any scalar multiple gives the same curve) and guarantees
	the discriminant condition $\Delta < 0$.
	The full $6 \times 6$ constraint matrix $\mathbf{C}$ is
	constructed so that $\mathbf{a}^\top\mathbf{C}\mathbf{a}$
	evaluates exactly to $4ac - b^2$:
	
	\begin{equation}
		\mathbf{C} =
		\begin{bmatrix}
			0 & 0 & 2 & 0 & 0 & 0 \\
			0 &-1 & 0 & 0 & 0 & 0 \\
			2 & 0 & 0 & 0 & 0 & 0 \\
			0 & 0 & 0 & 0 & 0 & 0 \\
			0 & 0 & 0 & 0 & 0 & 0 \\
			0 & 0 & 0 & 0 & 0 & 0
		\end{bmatrix}
	\end{equation}
	
	The lower $3\times3$ block is zero because $d$, $e$, $f$
	do not appear in the discriminant condition.
	
	
	\section{Design Matrix and Scatter Matrix}
	
	
	\textbf{Goal:} Organise all $n$ data points into a compact
	matrix form that the optimisation algorithm can work with.
	. Each point $(x_k, y_k)$ contributes one row
	to the design matrix $\mathbf{D}$ by evaluating the six
	conic monomials.
	Multiplying $\mathbf{D}^\top\mathbf{D}$ accumulates
	information from all $n$ points into the fixed-size
	$6\times6$ scatter matrix $\mathbf{S}$.
	
	\subsection{Design Matrix}
	
	Each point $(x_k, y_k)$ contributes one row:
	
	\begin{equation}
		\mathbf{d}_k = \bigl[x_k^2,\; x_k y_k,\; y_k^2,\;
		x_k,\; y_k,\; 1\bigr]
		\in \mathbb{R}^{1 \times 6}
	\end{equation}
	
	The full design matrix stacking all $n$ points is:
	
	\begin{equation}
		\mathbf{D} =
		\begin{bmatrix}
			x_1^2 & x_1 y_1 & y_1^2 & x_1 & y_1 & 1 \\
			x_2^2 & x_2 y_2 & y_2^2 & x_2 & y_2 & 1 \\
			\vdots & \vdots & \vdots & \vdots & \vdots & \vdots \\
			x_n^2 & x_n y_n & y_n^2 & x_n & y_n & 1
		\end{bmatrix}
		\in \mathbb{R}^{n \times 6}
	\end{equation}
	
	$\mathbf{D}$ is partitioned by column type
	(quadratic vs.\ linear monomials):
	
	\begin{equation}
		\mathbf{D} = [\,\mathbf{D}_1\; |\; \mathbf{D}_2\,],
		\qquad
		\mathbf{D}_1 = [x^2,\, xy,\, y^2],\quad
		\mathbf{D}_2 = [x,\, y,\, 1]
	\end{equation}
	
	with corresponding coefficient partition
	$\mathbf{a}_1 = [a,b,c]^\top$ and $\mathbf{a}_2 = [d,e,f]^\top$.
	
	\subsection{Scatter Matrix}
	
	The scatter matrix accumulates all $n$ points
	into a $6\times6$ summary via:
	
	\begin{equation}
		\mathbf{S} = \mathbf{D}^\top \mathbf{D} \in \mathbb{R}^{6\times6},
		\qquad S_{ij} = \sum_{k=1}^{n} D_{ki}\, D_{kj}
	\end{equation}
	
	Every entry is a sum over all $n$ points.
	For example $S_{11} = \sum x_k^4$ and $S_{66} = n$.
	No information is lost: $\mathbf{S}$ is a complete
	statistical summary of the point cloud.
	Partitioned into $3\times3$ blocks:
	
	\begin{equation}
		\mathbf{S} =
		\begin{bmatrix}
			\mathbf{S}_{11} & \mathbf{S}_{12} \\
			\mathbf{S}_{21} & \mathbf{S}_{22}
		\end{bmatrix},
		\qquad
		\mathbf{S}_{21} = \mathbf{S}_{12}^\top
	\end{equation}
	
	where explicitly:
	
	\begin{align}
		\mathbf{S}_{11} &=
		\begin{bmatrix}
			\sum x^4     & \sum x^3 y   & \sum x^2 y^2 \\
			\sum x^3 y   & \sum x^2 y^2 & \sum x y^3   \\
			\sum x^2 y^2 & \sum x y^3   & \sum y^4
		\end{bmatrix}, \quad
		\mathbf{S}_{22} =
		\begin{bmatrix}
			\sum x^2 & \sum x y & \sum x \\
			\sum x y & \sum y^2 & \sum y \\
			\sum x   & \sum y   & n
		\end{bmatrix}
	\end{align}
	
	
	\section{Constrained Least-Squares Problem}

	
	\textbf{Goal:} Find the coefficient vector $\mathbf{a}$
	that makes the ellipse pass as close as possible to all
	$n$ data points while remaining a valid ellipse.
	Minimizing the total squared algebraic distance
	$\mathbf{a}^\top\mathbf{S}\mathbf{a}$ subject to the
	ellipse constraint, using a Lagrange multiplier.
	This converts the constrained optimisation into a
	generalised eigenvalue problem.
	
	\subsection{Lagrangian Formulation}
	
	We minimise the sum of squared algebraic distances subject to
	the ellipse constraint \eqref{eq:constraint}:
	
	\begin{equation}
		\min_{\mathbf{a}}\; \mathbf{a}^\top \mathbf{S}\, \mathbf{a}
		\quad \text{subject to} \quad
		\mathbf{a}^\top \mathbf{C}\, \mathbf{a} = 1
	\end{equation}
	
	Introducing Lagrange multiplier $\lambda$ and differentiating:
	
	\begin{equation}
		\mathcal{L}(\mathbf{a}, \lambda) =
		\mathbf{a}^\top \mathbf{S}\, \mathbf{a}
		- \lambda\bigl(\mathbf{a}^\top \mathbf{C}\, \mathbf{a} - 1\bigr)
	\end{equation}
	
	Setting $\partial \mathcal{L}/\partial \mathbf{a} = \mathbf{0}$
	gives the \textbf{generalised eigenvalue problem}:
	
	\begin{equation}
		\boxed{\mathbf{S}\, \mathbf{a} = \lambda\, \mathbf{C}\, \mathbf{a}}
		\label{eq:geneig}
	\end{equation}
	
	Pre-multiplying by $\mathbf{a}^\top$ shows that
	$\lambda = \mathbf{a}^\top\mathbf{S}\mathbf{a}$, i.e.\ $\lambda$
	equals the objective (fitting error).
	The smallest $|\lambda|$ gives the best fit.
	

	\section{Schur Complement Reduction ($6\times6 \to 3\times3$)}
	
	
	\textbf{Goal:} Reduce the expensive $6\times6$ generalised
	eigenvalue problem to a cheap $3\times3$ standard one.
	Exploiting the block structure of $\mathbf{C}$
	(its lower $3\times3$ block is zero) to eliminate $\mathbf{a}_2$
	algebraically from the system, leaving only a $3\times3$
	problem in $\mathbf{a}_1$.
	
	\subsection{Block System}
	
	Writing \eqref{eq:geneig} in block form:
	
	\begin{equation}
		\begin{bmatrix}
			\mathbf{S}_{11} & \mathbf{S}_{12} \\
			\mathbf{S}_{21} & \mathbf{S}_{22}
		\end{bmatrix}
		\begin{bmatrix} \mathbf{a}_1 \\ \mathbf{a}_2 \end{bmatrix}
		= \lambda
		\begin{bmatrix}
			\mathbf{C}_1 & \mathbf{0} \\
			\mathbf{0}   & \mathbf{0}
		\end{bmatrix}
		\begin{bmatrix} \mathbf{a}_1 \\ \mathbf{a}_2 \end{bmatrix}
	\end{equation}
	
	where $\mathbf{C}_1$ is the top-left $3\times3$ block:
	
	\begin{equation}
		\mathbf{C}_1 =
		\begin{bmatrix}
			0 & 0 & 2 \\
			0 &-1 & 0 \\
			2 & 0 & 0
		\end{bmatrix}
	\end{equation}
	
	This gives two block equations:
	
	\begin{align}
		\mathbf{S}_{11}\,\mathbf{a}_1 + \mathbf{S}_{12}\,\mathbf{a}_2
		&= \lambda\, \mathbf{C}_1\,\mathbf{a}_1
		\tag{I}\label{eq:blockI} \\[4pt]
		\mathbf{S}_{21}\,\mathbf{a}_1 + \mathbf{S}_{22}\,\mathbf{a}_2
		&= \mathbf{0}
		\tag{II}\label{eq:blockII}
	\end{align}
	
	\subsection{Elimination of $\mathbf{a}_2$}
	
	Equation (II) contains no $\lambda$ (because the lower block
	of $\mathbf{C}$ is zero), so it can be solved directly for $\mathbf{a}_2$:
	
	\begin{equation}
		\mathbf{a}_2 = -\mathbf{S}_{22}^{-1}\, \mathbf{S}_{12}^\top\, \mathbf{a}_1
		\label{eq:backsubst}
	\end{equation}
	
	This is the back-substitution formula: once $\mathbf{a}_1$
	is known, $\mathbf{a}_2$ is completely determined.
	
	\subsection{Schur Complement}
	
	Substituting \eqref{eq:backsubst} into \eqref{eq:blockI}
	and factoring out $\mathbf{a}_1$ yields:
	
	\begin{equation}
		\boxed{\mathbf{S}^* = \mathbf{S}_{11}
			- \mathbf{S}_{12}\,\mathbf{S}_{22}^{-1}\,\mathbf{S}_{12}^\top}
		\label{eq:schur}
	\end{equation}
	
	$\mathbf{S}^*$ is the \textbf{Schur complement}: the part of
	the quadratic information in $\mathbf{S}_{11}$ that cannot be
	explained by the linear terms.
	The $6\times6$ problem reduces to a $3\times3$ generalised problem:
	
	\begin{equation}
		\mathbf{S}^*\,\mathbf{a}_1 = \lambda\,\mathbf{C}_1\,\mathbf{a}_1
	\end{equation}
	
	Pre-multiplying by $\mathbf{C}_1^{-1}$ converts to standard form:
	
	\begin{equation}
		\underbrace{\mathbf{C}_1^{-1}\,\mathbf{S}^*}_{\mathbf{A}}\,\mathbf{a}_1
		= \lambda\,\mathbf{a}_1
		\label{eq:standard_eig}
	\end{equation}
	

	\section{Analytical Inverse of $\mathbf{S}_{22}$}

	
	\textbf{Goal:} Compute $\mathbf{S}_{22}^{-1}$ needed for both
	the Schur complement and the back-substitution.
    Since $\mathbf{S}_{22}$ is always $3\times3$,
	its inverse is computed in closed form using the cofactor method ---
	no iterative algorithm needed.
	A small regularisation $\varepsilon = 10^{-6}$ is added to the
	diagonal first to prevent singularity when the data is
	nearly collinear.
	
	With regularisation $\varepsilon = 10^{-6}$ added to the diagonal, let:
	
	\begin{equation}
		\mathbf{S}_{22} + \varepsilon\mathbf{I} =
		\begin{bmatrix}
			\alpha & \beta  & \gamma \\
			\beta  & \delta & \zeta  \\
			\gamma & \zeta  & \eta
		\end{bmatrix}
	\end{equation}
	
	The determinant (required for all cofactor entries) is:
	
	\begin{equation}
		\det = \alpha(\delta\eta - \zeta^2)
		- \beta(\beta\eta - \zeta\gamma)
		+ \gamma(\beta\zeta - \delta\gamma)
	\end{equation}
	
	The inverse entries (exploiting symmetry $Q_{ij} = Q_{ji}$) are:
	
	\begin{align}
		Q_{11} &= \phantom{-}\frac{\delta\eta - \zeta^2}{\det} &
		Q_{12} &= -\frac{\beta\eta - \gamma\zeta}{\det} &
		Q_{13} &= \phantom{-}\frac{\beta\zeta - \gamma\delta}{\det} \\[6pt]
		Q_{22} &= \phantom{-}\frac{\alpha\eta - \gamma^2}{\det} &
		Q_{23} &= -\frac{\alpha\zeta - \gamma\beta}{\det} &
		Q_{33} &= \phantom{-}\frac{\alpha\delta - \beta^2}{\det}
	\end{align}
	

	\section{Matrix $\mathbf{A}$ and Eigenvalue Problem}
	
	
	\textbf{Goal:} Form the $3\times3$ matrix $\mathbf{A}$ and
	solve its eigenvalue problem to find the quadratic coefficients
	$\mathbf{a}_1 = [a,b,c]^\top$.
	Since $\mathbf{C}_1$ is a fixed constant matrix,
	its inverse $\mathbf{C}_1^{-1}$ is precomputed analytically
	and hardcoded.
	The eigenvector with the smallest $|\lambda|$ is selected
	because $\lambda$ equals the fitting error.
	
	\subsection{Computing $\mathbf{A}$}
	
	$\mathbf{A} = \mathbf{C}_1^{-1}\,\mathbf{S}^*$, where
	$\mathbf{C}_1^{-1}$ is a fixed constant hardcoded in the code
	(it never changes because $\mathbf{C}_1$ depends only on the
	mathematical ellipse constraint, not on any data):
	
	\begin{equation}
		\det(\mathbf{C}_1) = 4, \qquad
		\mathbf{C}_1^{-1} =
		\begin{bmatrix}
			0   & 0   & 0.5 \\
			0   & -1  & 0   \\
			0.5 & 0   & 0
		\end{bmatrix}
	\end{equation}
	
	\subsection{Standard Eigenvalue Problem}
	
	The $3\times3$ standard eigenvalue problem is solved:
	
	\begin{equation}
		\mathbf{A}\,\mathbf{a}_1 = \lambda\,\mathbf{a}_1
	\end{equation}
	
	yielding eigenvalue-eigenvector pairs
	$\{(\lambda_1, \mathbf{v}_1),(\lambda_2, \mathbf{v}_2),
	(\lambda_3, \mathbf{v}_3)\}$.
	The optimal quadratic coefficients are the eigenvector
	with the smallest absolute eigenvalue (smallest fitting error):
	
	\begin{equation}
		\mathbf{a}_1 = \mathbf{v}_{k^*}, \qquad
		k^* = \operatorname*{arg\,min}_{k}\, |\lambda_k|
	\end{equation}
	

	\section{Back-Substitution for Linear Coefficients}
	
	
	\textbf{Goal:} Recover the linear coefficients $[d,e,f]^\top$
	from the quadratic ones $[a,b,c]^\top$.
	Substituting $\mathbf{a}_1$ into the exact
	algebraic formula \eqref{eq:backsubst} derived in
	Section~6.2.
	This is not an approximation --- it is the unique
	linear-coefficient vector that minimises the fitting error
	given $\mathbf{a}_1$.
	
	Once $\mathbf{a}_1 = [a, b, c]^\top$ is known, the linear
	coefficients $\mathbf{a}_2 = [d, e, f]^\top$ follow exactly:
	
	\begin{equation}
		\mathbf{a}_2 = -\mathbf{S}_{22}^{-1}\,\mathbf{S}_{12}^\top\,\mathbf{a}_1
	\end{equation}
	
	The full six-parameter vector is assembled as:
	
	\begin{equation}
		\mathbf{a} =
		\begin{bmatrix} \mathbf{a}_1 \\ \mathbf{a}_2 \end{bmatrix}
		= [a,\; b,\; c,\; d,\; e,\; f]^\top
	\end{equation}
	

	\section{Fitzgibbon Normalisation}

	
	\textbf{Goal:} Fix the scale ambiguity of the eigenvector
	(which is only defined up to a scalar multiple) so that
	coefficients from different windows are comparable.
	Dividing $\mathbf{a}$ by $\sqrt{|4ac - b^2|}$,
	which forces the Fitzgibbon constraint to equal exactly 1.
	This does not change the ellipse shape --- only its
	numerical scale.
	
	The normalisation scale factor is:
	
	\begin{equation}
		\sigma = \sqrt{\bigl|\,\mathbf{a}^\top\mathbf{C}\,\mathbf{a}\,\bigr|}
		= \sqrt{|4ac - b^2|}
	\end{equation}
	
	\begin{equation}
		\mathbf{a}_{\text{norm}} = \frac{\mathbf{a}}{\sigma}
	\end{equation}
	
	\textbf{Verification} --- the constraint is satisfied after normalisation:
	
	\begin{equation}
		\mathbf{a}_{\text{norm}}^\top\,\mathbf{C}\,\mathbf{a}_{\text{norm}}
		= \frac{1}{\sigma^2}\,\mathbf{a}^\top\mathbf{C}\,\mathbf{a}
		= \frac{4ac - b^2}{|4ac - b^2|}
		= 1 \quad \checkmark
	\end{equation}
	
	
	\section{Conic to Parametric Conversion}

	
	\textbf{Goal:} Convert the six algebraic coefficients
	$[a,b,c,d,e,f]$ into human-interpretable geometric parameters:
	ellipse centre $(x_0, y_0)$, semi-axes $(r_1, r_2)$,
	and tilt angle $\theta$.
	Working through five sub-steps --- find the
	critical point of $F$ (centre), evaluate $F$ there,
	diagonalise the quadratic part (angle), compute the
	eigenvalues of the quadratic form (axis scaling),
	and extract the semi-axes.
	
	\subsection{Step 1: Ellipse Centre}
	
	The centre is the unique point where both partial derivatives
	of $F(x,y)$ vanish simultaneously
	($\partial F/\partial x = \partial F/\partial y = 0$),
	solved by Cramer's rule with denominator $\Delta = b^2 - 4ac < 0$:
	
	\begin{equation}
		\begin{bmatrix} 2a & b \\ b & 2c \end{bmatrix}
		\begin{bmatrix} x_0 \\ y_0 \end{bmatrix}
		=
		\begin{bmatrix} -d \\ -e \end{bmatrix}
		\;\implies\;
		x_0 = \frac{2cd - be}{b^2 - 4ac}, \qquad
		y_0 = \frac{2ae - bd}{b^2 - 4ac}
		\label{eq:centre}
	\end{equation}
	
	\subsection{Step 2: Value at Centre}
	
	Evaluating $F$ at the centre gives the scalar $a'$,
	which acts as a normalisation constant for the semi-axes.
	For a valid (non-degenerate) ellipse, $a' < 0$:
	
	\begin{equation}
		a' = ax_0^2 + bx_0 y_0 + cy_0^2 + dx_0 + ey_0 + f
		\label{eq:aprime}
	\end{equation}
	
	\subsection{Step 3: Tilt Angle}
	
	The tilt angle $\theta$ is the rotation that eliminates the
	cross-term $bxy$ by diagonalising the $2\times2$ quadratic
	form matrix $\mathbf{Q} =
	\bigl[\begin{smallmatrix}a&b/2\\b/2&c\end{smallmatrix}\bigr]$:
	
	\begin{equation}
		\tan(2\theta) = \frac{b}{a - c}
		\;\implies\;
		\theta =
		\begin{cases}
			0       & \text{if } b \approx 0
			\;\text{(already axis-aligned)}\\
			\pi/4   & \text{if } a \approx c
			\;\text{(isotropic quadratic part)}\\
			\dfrac{1}{2}\arctan\!\left(\dfrac{b}{a-c}\right)
			& \text{otherwise}
		\end{cases}
		\label{eq:angle}
	\end{equation}
	
	\subsection{Step 4: Eigenvalues of the Quadratic Form}
	
	The eigenvalues of $\mathbf{Q}$ determine the curvature of
	the ellipse in each axis direction.
	They are found from the characteristic equation of the
	$2\times2$ matrix $\mathbf{Q}$:
	
	\begin{equation}
		\lambda_{1,2} = \frac{(a + c) \pm \sqrt{(a-c)^2 + b^2}}{2}
		\label{eq:lambdas}
	\end{equation}
	
	\subsection{Step 5: Semi-Axes}
	
	After translating to the centre and rotating by $\theta$,
	the conic reduces to the canonical form
	$\lambda_1 u^2 + \lambda_2 v^2 = -a'$.
	Comparing with the standard ellipse form
	$u^2/r_1^2 + v^2/r_2^2 = 1$ gives the semi-axes:
	
	\begin{equation}
		r_1 = \sqrt{\frac{-a'}{\lambda_1}}, \qquad
		r_2 = \sqrt{\frac{-a'}{\lambda_2}}
		\label{eq:semiaxes}
	\end{equation}
	
	with $r_{\text{major}} = \max(r_1, r_2)$ and
	$r_{\text{minor}} = \min(r_1, r_2)$.
	
	\subsection{Parametric Form}
	
	The complete fitted ellipse in the original coordinate system:
	
	\begin{equation}
		\begin{bmatrix} x(t) \\ y(t) \end{bmatrix}
		=
		\begin{bmatrix} x_0 \\ y_0 \end{bmatrix}
		+
		\begin{bmatrix}
			\cos\theta & -\sin\theta \\
			\sin\theta &  \cos\theta
		\end{bmatrix}
		\begin{bmatrix} r_1 \cos t \\ r_2 \sin t \end{bmatrix},
		\qquad t \in [0,\, 2\pi)
	\end{equation}
	
%
%	\section{Impact Onset Detection}
%	
%	
%	\textbf{Goal:} Automatically detect the exact time when the
%	impact begins in the raw signal.
%	\textbf{How:} Find the peak of the signal, then walk backward
%	in time until the signal last fell below the 2-sigma baseline
%	threshold.
%	This backward-search approach gives the true onset with
%	zero forward-window lag.
%	
%	\subsection{Stable Baseline Statistics}
%	
%	A stable pre-impact region is used as the reference.
%	Its mean and standard deviation are:
%	
%	\begin{equation}
%		\mu_{\text{bl}} = \frac{1}{N_{\text{bl}}}
%		\sum_{k \in \mathcal{B}} x(k),
%		\qquad
%		\sigma_{\text{bl}} = \sqrt{\frac{1}{N_{\text{bl}}}
%			\sum_{k \in \mathcal{B}} \bigl(x(k) - \mu_{\text{bl}}\bigr)^2}
%	\end{equation}
%	
%	\subsection{Backward Search for Onset}
%	
%	Let $k^* = \arg\max_k |x(k)|$ be the impact peak index.
%	Walking backward from $k^*$, the onset index $k_0$ is the
%	last sample before the peak where the signal was
%	still within the baseline level:
%	
%	\begin{equation}
%		k_0 = \max\bigl\{k \leq k^* :
%		|x(k) - \mu_{\text{bl}}| < 2\sigma_{\text{bl}}\bigr\} + 1
%	\end{equation}
%	
%	The impact onset time is $t_{\text{impact}} = k_0 \cdot \Delta t$.
%	
%
%	\section{Change Point Detection for Parameter $a$}
%	
%	
%	\textbf{Goal:} Detect all times when the ellipse coefficient
%	$a$ (blue curve) changes significantly from its pre-impact
%	baseline, marking each as a dashed vertical line in the plots.
%	\textbf{How:} Compute the stable baseline statistics of $a$
%	from the 0.5--3.0\,s window, flag all windows where $a$
%	deviates by more than $3\sigma$, and group nearby detections
%	separated by at least 0.3\,s.
%	
%	\subsection{Stable Baseline of Parameter $a$}
%	
%	Using the stable time window $t \in [0.5\,\text{s},\, 3.0\,\text{s}]$
%	(which excludes the startup transient and the impact region):
%	
%	\begin{equation}
%		\bar{a}_{\text{bl}} = \frac{1}{|\mathcal{S}|}
%		\sum_{i \in \mathcal{S}} a_i, \qquad
%		\sigma_{a}^{\text{bl}} = \sqrt{\frac{1}{|\mathcal{S}|}
%			\sum_{i \in \mathcal{S}} (a_i - \bar{a}_{\text{bl}})^2}
%	\end{equation}
%	
%	\subsection{Three-Sigma Change Detection}
%	
%	A statistically significant change is declared at window $i$
%	when $a_i$ deviates more than 3 standard deviations from
%	the baseline (probability $<0.3\%$ for Gaussian noise):
%	
%	\begin{equation}
%		|a_i - \bar{a}_{\text{bl}}| > 3\,\sigma_{a}^{\text{bl}}
%	\end{equation}
%	
%	The first sample of each new group of consecutive detections
%	(a rising edge in the binary detection signal) is recorded
%	as a change point, with groups separated by at least
%	$\Delta t_{\min} = 0.3$\,s.
	

	\section{Min-Max Normalisation for Stacked Visualisation}

	
	\textbf{Goal:} Display all five ellipse parameters
	($a$, $b$, $c$, $d$, $e$) on the same plot with equal
	visual amplitude so they can be compared directly.
	Normalising the Min-max of each parameter to $[0,1]$
	independently, then add a fixed vertical offset to
	separate the five bands.
	This guarantees equal band height regardless of the
	raw magnitude of each parameter.
	
	For each parameter $j \in \{1,\ldots,5\}$:
	
	\begin{equation}
		v_{\text{norm},i}^{(j)}
		= \frac{v_i^{(j)} - \min_i v_i^{(j)}}
		{\max_i v_i^{(j)} - \min_i v_i^{(j)}}
		\in [0,\, 1]
	\end{equation}
	
	The stacked (offset) display value is:
	
	\begin{equation}
		\tilde{v}_i^{(j)} = v_{\text{norm},i}^{(j)} + \delta_j, \qquad
		\delta_j = (5 - j) \times \Delta_{\text{sep}}
	\end{equation}
	
	where $\Delta_{\text{sep}} = 1.15$ is the vertical separation,
	placing parameter $a$ ($j=1$) at the top (offset $= 4\times1.15$)
	and $e$ ($j=5$) at the bottom (offset $= 0$).
	Y-axis tick labels are placed at the centre of each band:
	
	\begin{equation}
		y_{\text{tick}}^{(j)} = \delta_j + 0.5
	\end{equation}
	
\end{document}